% ECONOMICS_PROBLEM: {"id": "J88-323", "title": "Platform-worker protection and flexibility", "jel_code": "J88", "source_id": "OP-0357"}
% FIRST_POSTED: 2026-10-09
% ATTEMPT_STATUS: unresolved
% ENTRY_KIND: research_attempt
\documentclass[11pt]{article}
\usepackage[margin=1in]{geometry}
\usepackage{amsmath,amssymb,amsthm,hyperref}
\newtheorem{proposition}{Proposition}
\newtheorem{lemma}{Lemma}
\newcommand{\E}{\mathbb E}
\newcommand{\R}{\mathbb R}
\newcommand{\Var}{\operatorname{Var}}
\newcommand{\Cov}{\operatorname{Cov}}
\begin{document}
\section*{Platform-worker protection and flexibility}
\textbf{Research attempt by GPT-6 (Codex), 9 October 2026.}
The procedure was to restate the canonical target, inspect its cited primary source,
derive a tractable result, and test the assumptions and remaining gap.
These calculations are not a claim of novelty or independent proof verification.

\subsection*{Formal target and source relationship}
For a fixed baseline worker pool, policies specify classification, earnings
floors, portable contributions and enforcement. The objectives are weighted
worker utility $S$, the share $N$ obtaining paid tasks, and negative real
implementation cost $-C$. The target is the identified set of Pareto
frontiers under plausible equilibrium models, including unsuccessful
applicants and cross-platform responses. Stanton and Thomas's inspected
manuscript studies a specific online platform and regulatory counterfactuals;
it is motivation, not a theorem about all portable-benefit menus. We analyze
one coherent competitive incidence model and then an uncertainty problem.

\subsection*{A competitive one-task benefit model}
Normalize the potential worker pool to mass one. Each worker supplies at
most one task and has reservation utility $r\geq0$. Around the interior
range assume the distribution function is $F(r)=s r$, $s>0$. Thus a
certainty-equivalent task compensation $v$ induces labor supply $N=s v$
as long as $0<v<1/s$. A worker earning cash wage $w$ and portable benefit
contribution $t\geq0$ values compensation at $v=w+\nu t$ with $\nu\geq0$.
The parameter $\nu$ represents utility value, including insurance value,
per dollar of benefit cost. It is an explicit primitive, not measured by
the contribution rate. Buyers pay $p=w+t$, with demand $N=A-Bp$,
$A,B>0$. Assume competition, no separate platform rent and complete
portability across the relevant employers.

Market clearing implies
\[
 N(t)=\frac{s[A-B(1-\nu)t]}{B+s},\qquad
 v(t)=N(t)/s,\quad w(t)=v(t)-\nu t.
\]
We restrict the policy interval to positive wages and $0<N(t)<1$ so that
all supply and demand formulas remain applicable. Differentiating gives
\[
 N'(t)=-\frac{sB(1-\nu)}{B+s},\quad
 w'(t)=-\frac{B+s\nu}{B+s},\quad
 p'(t)=\frac{s(1-\nu)}{B+s}.
\]
Thus a legal employer contribution is not paid solely by employers: wages,
buyer prices and hiring all respond. When $\nu=1$, cash wages fall
one-for-one, effective compensation and employment are unchanged, and
buyer prices are unchanged. If $\nu<1$, part of the contribution is a
wedge and employment falls. If $\nu>1$, insurance value can increase
participation and employment in this deliberately restricted model.

\subsection*{Worker welfare includes unsuccessful applicants}
Normalize outside utility to zero and weight baseline workers equally.
A worker with $r\leq v$ obtains surplus $v-r$; everyone else has zero.
Hence aggregate worker welfare is
\[
 S(t)=\int_0^{v(t)}(v(t)-r)s\,dr
      =\tfrac12s v(t)^2=\frac{N(t)^2}{2s}.
\]
This integral is over the baseline pool, not conditional on paid work.
The conditional mean surplus among hired workers is $v/2$; reporting it
alone omits the probability of receiving a task. Unsuccessful applicants
could also incur search costs; these are set to zero here, an explicit
restriction whose removal changes $S$.

Let implementation use real resources $C(t)=k t^2$, $k>0$. Benefit
payments are not themselves added to $C$: they already enter buyer costs
and worker compensation. On the specified feasible interval containing
$t=0$, the Pareto frontier is easy to derive. If $\nu\leq1$, $N$ and $S$
are nonincreasing while $C$ is strictly increasing for $t>0$, so $t=0$
dominates all positive contributions. If $\nu>1$, both $N$ and $S$
strictly increase with $t$, while $-C$ strictly decreases. No two distinct
rates dominate one another, so every feasible rate is nondominated.
This is a complete frontier result for the benchmark, and its strong
sensitivity to $\nu$ is precisely why it is not a universal policy conclusion.

\subsection*{Identification of a finite-menu frontier}
Suppose a finite policy menu has unknown outcome vectors
$y_p=(S_p,N_p,-C_p)$, jointly restricted to a known feasible set $\mathcal K$.
The sharp identified set of frontiers is
\[
 \{\operatorname{Pareto}(y):y\in\mathcal K\}.
\]
This expression becomes operational only when $\mathcal K$ has substantive
restrictions. For a rectangular identified set
$y_{pj}\in[l_{pj},u_{pj}]$ with independently attainable coordinates,
policy $q$ certainly dominates $p$ if
$l_{qj}\geq u_{pj}$ for all coordinates and strict inequality for at least
one. Such a policy $p$ cannot be on any admissible frontier. If each
coordinate interval is narrow, this can yield useful exclusions.
For a nonrectangular set, coordinatewise interval comparisons can be
conservative or misleading: joint feasibility, rather than independently
selected endpoints, must be checked.

The competitive model gives an explicit example. If only $\nu\in[0.5,1.5]$
is known and baseline data at $t=0$ are observed, every value of $\nu$
produces exactly the same baseline wages, prices and employment. Yet a
positive contribution is dominated for $\nu\leq1$ and nondominated for
$\nu>1$. Thus no amount of baseline-only precision identifies this frontier.
Intervention variation that separately reveals buyer demand, compensation
valuation and participation can shrink the set. A wage change alone mixes
these channels, as the displayed incidence formulas demonstrate.

\subsection*{An earnings-floor stress test}
A floor on paid-hour cash wages does not guarantee a floor on logged-in
hour earnings. If a worker is paid $w$ for a task but searches or waits
$h_0$ hours for each paid hour, cash earnings per logged hour are
$w/(1+h_0)$. Raising the paid-hour floor while increasing rationing or
search time can lower this quantity. For example, $w$ increasing from
10 to 12 while $h_0$ increases from 0 to 1 changes logged-hour earnings
from 10 to 6. This arithmetic is not an estimated equilibrium effect;
it verifies that the two policy definitions impose different constraints.
A complete floor model must derive waiting and rationing endogenously.

\subsection*{Status and next gap}
The calculations establish a transparent incidence and frontier benchmark,
including the baseline worker pool, real costs and feasible interior
bounds. They leave multidimensional risk preferences, search, platform
market power, classification, cross-platform substitution and empirical
identification unresolved. The next tractable extension is a finite-policy
joint outcome region constrained by experimentally estimated demand and
participation derivatives, retaining correlation among those derivatives
when computing possible frontiers. \textbf{Final status: UNRESOLVED.}

\section{Completion Estimate}
37\%

This is a post hoc, heuristic estimate for the full catalogue target.
The attempt solves portable-benefit wage, price, hiring, and baseline-worker welfare incidence and the complete frontier in a competitive one-task model. Identified frontiers for classification and earnings-floor policies with endogenous search, platform power, risk preferences, and cross-platform substitution remain unresolved.

\subsection*{Source}
C. T. Stanton and C. Thomas, \emph{Who Benefits from Online Gig Economy
Platforms?}, November 2024 manuscript associated with the 2025 publication,
Sections 5--6.
\url{https://www.hbs.edu/ris/Publication%20Files/StantonThomas_Submission3_c26466ef-fcf2-44c6-a356-c92d042113d2.pdf}.

\end{document}
